\chapter{Ciblage de volatilité et portefeuilles gérés du risque}\label{ch:s1:volatility-targeting-and-risk-managed-portfolios}

Lorsque la volatilité grimpe, chaque portefeuille dimensionné sur une cible de volatilité doit vendre, le même jour et pour la même raison. Moreira et Muir ont montré
que prendre moins de risque lorsque la volatilité est élevée a relevé les ratios de Sharpe du marché et de la plupart des facteurs actions~; Cederburg, O'Doherty, Wang et Yan,
en testant 103 stratégies comme le ferait un investisseur en temps réel, ont constaté que les \omterm{def:s1:volatility-targeting-and-risk-managed-portfolios:managed}{portefeuilles pilotés par la volatilité} ne battent pas systématiquement les portefeuilles non gérés. Sur
les marchés synthétiques de ce chapitre, les deux affirmations sont vraies~: piloter le marché par sa volatilité améliore le ratio de Sharpe de 0.03 en moyenne sur vingt
échantillons de dix ans, et ne fait mieux que dans la moitié d'entre eux. Ce qui est certain, c'est la vente~: après une chute synthétique de 6~\%, les fonds détenant un cinquième d'un pour cent
du marché vendent 14.5~\% du volume du lendemain. La construction repose sur \texttt{firm.voltarget}.

\section{Mécanique du ciblage de volatilité}

Le chapitre~19 a défini le \omterm{def:s1:trend-following:voltarget}{ciblage de volatilité}~: une position dimensionnée à une volatilité cible divisée par sa prévision. Un portefeuille géré du risque ajoute une règle
pour la prévision et pour la fréquence des transactions.

\begin{definition}[Portefeuille piloté par la volatilité]\label{def:s1:volatility-targeting-and-risk-managed-portfolios:managed}
Un \emph{portefeuille piloté par la volatilité}\index{portefeuille piloté par la volatilité} ajuste l'exposition d'un portefeuille sous-jacent par l'inverse de sa volatilité prévue
(ou de sa variance), de sorte qu'il prend moins de risque après les périodes agitées et davantage après les périodes calmes, généralement avec un plafond sur l'effet de levier et une bande
à l'intérieur de laquelle l'exposition n'est pas rééquilibrée.
\end{definition}

Les ingrédients sont peu nombreux (\cref{lst:s1:volatility-targeting-and-risk-managed-portfolios:weights})~: une prévision de variance (la variance réalisée du mois précédent,
ou une variance pondérée exponentiellement comme dans le livre~4, chapitre~18)~; une règle de mise à l'échelle, cible divisée par la volatilité, ou cible au carré divisée par la variance selon Moreira et Muir, qui fait varier l'exposition plus vite~;
un plafond d'effet de levier~; une bande de rééquilibrage. Chaque règle de prévision et chaque bande produit un taux de rotation différent, et ce taux
constitue le coût de la stratégie.

\section{Améliore-t-il le ratio de Sharpe~?}

L'argument de Moreira et Muir est simple~: si les rendements espérés ne croissent pas proportionnellement à la volatilité, réduire l'exposition lorsque la volatilité est élevée
retire du risque plus vite qu'il ne retire du rendement. Le facteur de marché de \texttt{firm.synthmkt} satisfait cette hypothèse par construction~: une prime constante de 6~\% et une volatilité GJR-GARCH
avec effet de levier. Vingt échantillons de dix ans de ce facteur, détenus à exposition constante ou pilotés par la variance réalisée du mois précédent
(cible de 16~\%, effet de levier plafonné à trois)~:

\begin{center}
\small
\begin{tabular}{@{}p{5.6cm}>{\centering\arraybackslash}p{2.9cm}>{\centering\arraybackslash}p{2.9cm}>{\centering\arraybackslash}p{2.9cm}@{}}
\toprule
vingt échantillons de dix ans & exposition constante & inverse de la volatilité & inverse de la variance\\
\midrule
ratio de Sharpe moyen & 0.49 & 0.52 & 0.48\\
échantillons meilleurs que la constante & --- & 10 & 8\\
drawdown maximal moyen (volatilités annuelles) & 2.70 & 2.58 & 2.67\\
\bottomrule
\end{tabular}
\end{center}

Le gain moyen du portefeuille à inverse de la volatilité est de 0.028, avec un écart type de 0.22 entre échantillons
(\cref{fig:s1:volatility-targeting-and-risk-managed-portfolios:seeds}). Le ratio de Sharpe réel du portefeuille constant vaut $0.06/0.16 = 0.375$~; les
échantillons vont de $-0.44$ à 1.18. Un gain de quelques centièmes est réel dans le modèle et invisible sur une décennie quelconque, et la variance inverse, plus agressive,
le perd dans le bruit de ses propres prévisions. C'est le constat de Cederburg et de ses coauteurs en laboratoire~: l'amélioration existe dans la
population, et un investisseur en temps réel qui ne dispose que d'un seul historique ne peut pas compter dessus.

\begin{omfigure}
\begin{tikzpicture}
\begin{axis}[width=8cm,height=7cm,xlabel={\small ratio de Sharpe, exposition constante},ylabel={\small ratio de Sharpe, piloté par la volatilité},
  tick label style={font=\footnotesize},xmin=-0.6,xmax=1.4,ymin=-0.6,ymax=1.4,grid=major,grid style={gray!25},table/col sep=comma]
\addplot[only marks,mark=*,mark size=1.8pt,omThm] table[x=constant,y=managed]{figdata/strategies-1/26-volatility-targeting-and-risk-managed-portfolios/seeds.csv};
\addplot[black!50,dashed] coordinates {(-0.6,-0.6) (1.4,1.4)};
\end{axis}
\end{tikzpicture}

\omcaption{Vingt échantillons de dix ans du facteur de marché synthétique (GJR-GARCH, prime constante)~: le ratio de Sharpe du marché piloté par la volatilité
(inverse de la volatilité réalisée du mois précédent, cible de 16~\%, effet de levier plafonné à trois) en fonction de celui de l'exposition constante~; les points au-dessus de la
ligne pointillée sont les échantillons où le pilotage a aidé. Données~: \texttt{s1\_voltarget.seeds}.}\label{fig:s1:volatility-targeting-and-risk-managed-portfolios:seeds}
\end{omfigure}

Harvey, Hoyle, Korgaonkar, Rattray, Sargaison et van Hemert ont localisé le gain~: il se trouve uniquement dans les actifs risqués comme les actions et le crédit, où
la volatilité augmente quand les cours baissent (l'effet de levier)~; pour les obligations, les devises et les matières premières, le ratio de Sharpe change à peine. Ce que fait le ciblage
partout, c'est réduire les queues de distribution, parce que les pires journées surviennent en général lorsque la volatilité est déjà élevée et que l'exposition est déjà faible.
Les classes de \texttt{firm.synthfut} présentent un regroupement de volatilité mais pas d'effet de levier, et les cibler à 10~\% modifie leurs ratios de Sharpe dans les deux sens~:
0.19 à 0.19 pour les actions, 0.35 à 0.28 pour les obligations, $-0.18$ à $-0.08$ pour les devises, $-0.21$ à $-0.25$ pour les matières premières. Leurs
drawdowns maximaux baissent pour les obligations (de 6.19 à 4.99 volatilités annuelles) et les devises (de 13.28 à 11.45) et augmentent pour les actions et les matières premières.

\section{L'empreinte de l'industrie}

\begin{definition}[Flux de rebalancement des produits à effet de levier]\label{def:s1:volatility-targeting-and-risk-managed-portfolios:flow}
Un \emph{flux de rebalancement des produits à effet de levier}\index{flux de rebalancement des produits à effet de levier} est le solde des transactions de portefeuilles qui suivent des règles mécaniques d'exposition
(cibles de volatilité, remises à zéro quotidiennes des fonds à effet de levier, parité de risque) en réaction aux variations des cours et de la volatilité~; comme beaucoup suivent des règles similaires, le
flux est concentré et prévisible à partir de ces règles et des mouvements récents du marché.
\end{definition}

Le simulateur de flux de \texttt{firm.voltarget} (\cref{lst:s1:volatility-targeting-and-risk-managed-portfolios:flow}) place des fonds détenant 0.2~\% de la valeur d'un marché
dans un marché qui échange chaque jour 0.4~\% de sa valeur. Ils visent une volatilité de 10~\% avec une prévision pondérée exponentiellement (centre de masse de 20 jours,
effet de levier plafonné à deux). Le marché est calme à 12~\% pendant 250 jours, chute de 6~\% en une journée et évolue à 40~\% pendant vingt jours. Les fonds démarrent avec une exposition
de 0.73~; le lendemain de la chute, leur prévision bondit et ils vendent 14.5~\% du volume du jour, et sur les vingt jours un total de 20.5~\% du volume d'une journée,
pour finir avec une exposition de 0.25 (\cref{fig:s1:volatility-targeting-and-risk-managed-portfolios:spike}). Tous les paramètres du flux sont des hypothèses~:
la taille des sommes investies en \omterm{def:s1:trend-following:voltarget}{ciblage de volatilité} n'est mesurée par aucune source consultée pour ce chapitre.

\begin{omfigure}
\begin{tikzpicture}
\begin{axis}[width=11cm,height=5cm,xlabel={\small jours depuis la chute de 6~\%},ylabel={\small flux (\% du volume du jour)},
  tick label style={font=\footnotesize},xmin=-10,xmax=49,ymin=-16,ymax=2,grid=major,grid style={gray!25},axis y line*=left,
  table/col sep=comma]
\addplot[ybar,bar width=3pt,fill=omThm!60,draw=omThm] table[x=day,y=flow_pct]{figdata/strategies-1/26-volatility-targeting-and-risk-managed-portfolios/spike.csv};
\end{axis}
\begin{axis}[width=11cm,height=5cm,xmin=-10,xmax=49,ymin=0,ymax=1,axis y line*=right,axis x line=none,ylabel={\small exposition},
  tick label style={font=\footnotesize},table/col sep=comma]
\addplot[omDef,thick] table[x=day,y=exposure]{figdata/strategies-1/26-volatility-targeting-and-risk-managed-portfolios/spike.csv};
\end{axis}
\end{tikzpicture}

\omcaption{Fonds de \omterm{def:s1:trend-following:voltarget}{ciblage de volatilité} autour d'une chute synthétique de 6~\%~: leur flux net en part du volume de chaque jour (barres, à gauche) et leur exposition
(ligne, à droite). Les fonds détiennent 0.2~\% de la valeur du marché, et le marché échange 0.4~\% de sa valeur par jour (hypothèse). Données~:
\texttt{s1\_voltarget.spike}.}\label{fig:s1:volatility-targeting-and-risk-managed-portfolios:spike}
\end{omfigure}

Le flux déplace les cours, et ces mouvements alimentent les prévisions. Avec un impact linéaire de 0.1 fois la part du flux dans le volume, la chute sur vingt jours s'approfondit
de 16.6~\% à 18.7~\%~; avec 0.3, jusqu'à 22.0~\%. La vente est prévisible à partir d'informations publiques (la volatilité propre du marché), ce qui en fait
une fiche de stratégie pour quelqu'un d'autre~: quiconque est fournisseur de liquidité auprès des fonds en désendettement est rémunéré pour cela.

\section{Fiches de stratégie}

\begin{strategyfile}{Portefeuille volatilité gérée}\label{strat:s1:volatility-targeting-and-risk-managed-portfolios:equity}
\sfield{Qui rémunère, et pourquoi} Personne en particulier~: les rendements espérés n'augmentent pas autant que la volatilité, donc réduire l'exposition dans les périodes volatiles améliore le
compromis rendement-risque, si c'est bien le cas.
\sfield{Instruments et plateformes} Contrats à terme sur indices actions~; portefeuilles de facteurs actions.
\sfield{Signal} Volatilité prévue (réalisée ou EWMA).
\sfield{Dimensionnement et exécution} Exposition cible rapportée à la prévision, plafonnée~; bandes pour limiter le taux de rotation.
\sfield{Coûts} Rotation issue des changements de prévision, la plus élevée avec la variance inverse.
\sfield{Comment elle meurt} Bruit de prévision~; instabilité de la relation entre volatilité et rendements.
\sfield{Horizon, capacité, infrastructure} Rééquilibrage mensuel ou quotidien~; grande capacité.
\sfield{Backtest honnête} Prévisions en temps réel et constantes de mise à l'échelle~; comparer directement au portefeuille non piloté, et non seulement par l'alpha.
\sfield{Sources} Moreira et Muir (2017)~; Cederburg, O'Doherty, Wang et Yan (2020)~; ce chapitre~: $+0.03$ en moyenne, meilleur dans 10 échantillons sur 20.
\end{strategyfile}

\begin{strategyfile}{Multi-actifs à volatilité ciblée}\label{strat:s1:volatility-targeting-and-risk-managed-portfolios:multi}
\sfield{Qui rémunère, et pourquoi} Les primes des classes d'actifs, à risque plus stable.
\sfield{Instruments et plateformes} Contrats à terme sur actions, obligations, crédit, matières premières.
\sfield{Signal} La volatilité prévue de chaque actif et du portefeuille.
\sfield{Dimensionnement et exécution} Mise à l'échelle au niveau de chaque actif et du portefeuille~; plafond d'effet de levier.
\sfield{Coûts} Modérés.
\sfield{Comment elle meurt} Corrélations qui sautent avec la volatilité, de sorte que la diversification disparaît lorsque la cible réduit l'exposition.
\sfield{Horizon, capacité, infrastructure} De quotidien à mensuel~; un système de gestion des risques.
\sfield{Backtest honnête} Prévisions fondées uniquement sur des données passées~; statistiques de queue ainsi que ratios de Sharpe.
\sfield{Sources} Harvey, Hoyle, Korgaonkar, Rattray, Sargaison et van Hemert (2018).
\end{strategyfile}

\begin{strategyfile}{Trading contre le flux de désendettement}\label{strat:s1:volatility-targeting-and-risk-managed-portfolios:against}
\sfield{Qui rémunère, et pourquoi} Les fonds de \omterm{def:s1:trend-following:voltarget}{ciblage de volatilité} et autres fonds mécaniques, qui doivent vendre après les sauts de volatilité.
\sfield{Instruments et plateformes} Contrats à terme sur indices, à proximité de la clôture, lorsque les fonds rééquilibrent.
\sfield{Signal} La vente prévue~: les règles des fonds appliquées aux rendements récents du marché et une estimation de leur taille.
\sfield{Dimensionnement et exécution} Acheter dans le sens de la vente prévue, sortir à mesure que la pression s'atténue~; position modeste face au risque de krach.
\sfield{Coûts} Faibles~; le risque est la poursuite de la baisse.
\sfield{Comment elle meurt} Baisses qui se poursuivent pour des raisons fondamentales~; concurrence d'autres acteurs qui lisent les mêmes flux.
\sfield{Horizon, capacité, infrastructure} Quelques jours~; une estimation des actifs et des règles des fonds.
\sfield{Backtest honnête} Seulement l'information disponible à la date~; la poursuite du krach dans l'échantillon.
\sfield{Sources} Aucun chiffre de performance vérifié~; modèle de flux de ce chapitre.
\end{strategyfile}

\section{Tutoriel~: vente mécanique}\label{tut:s1:volatility-targeting-and-risk-managed-portfolios:tut}

\begin{tutorial}
\textbf{Objectif.} Gérer la volatilité du marché synthétique sur vingt échantillons et quatre classes d'actifs, et simuler la vente des fonds de ciblage après un
pic, avec et sans rétroaction. \textbf{État final~:} le tableau et les deux figures.
\begin{tutsteps}
\item \textbf{Prévisions et pondérations}.
\omcode{code/firm/voltarget/firm_voltarget.py}{26}{48}{Prévisions de variance et pondérations cibles.}\label{lst:s1:volatility-targeting-and-risk-managed-portfolios:weights}
\item \textbf{Le flux agrégé}~: les fonds de \omterm{def:s1:trend-following:voltarget}{ciblage de volatilité}, leur transaction en part du volume, et son impact réinjecté.
\omcode{code/firm/voltarget/firm_voltarget.py}{70}{91}{Flux des fonds de ciblage et sa rétroaction.}\label{lst:s1:volatility-targeting-and-risk-managed-portfolios:flow}
\item \textbf{Exécuter} \texttt{summary()}, \texttt{classes()} et \texttt{spike(impact)} pour trois impacts, et \texttt{fig\_voltarget.py}.
\end{tutsteps}
\textbf{Pour aller plus loin.} Ajouter une bande de rééquilibrage et mesurer son effet sur la rotation et sur le gain de Sharpe~; laisser plusieurs types de fonds (à \omterm{def:s1:trend-following:voltarget}{ciblage de volatilité},
remises à zéro quotidiennes des ETF à effet de levier) négocier ensemble~; laisser la taille des fonds croître au fil du temps.
\end{tutorial}

\section{Construction~: ciblage de volatilité}\label{bld:s1:volatility-targeting-and-risk-managed-portfolios:voltarget}

\begin{build}
\bfield{Rôle} Prévisions de variance, pondérations cibles avec plafonds et bandes, P\&L avec coûts, et simulateur de flux agrégé avec rétroaction.
\bfield{Interface} \texttt{realised\_var(r, window)}, \texttt{ewma\_var(r, com)}, \texttt{weights(var, target, cap, power)}, \texttt{banded(w, band)},
\texttt{run(r, w, cost)}, \texttt{spike\_flows(base, share, turnover, impact, target, com, cap)}.
\bfield{Règles} Prévisions fondées sur les rendements jusqu'à la clôture précédente~; pondérations plafonnées~; flux négociés le jour où ils sont décidés et réinjectés dans la prévision.
\bfield{Tests d'acceptation} \texttt{code/firm/voltarget/tests/}~: prévisions sur une série constante, pondérations, la bande et le P\&L à la main, vente après une baisse
et une baisse plus profonde avec rétroaction.
\bfield{Pour aller plus loin} Plusieurs types de fonds~; bandes~; calendrier intrajournalier des flux.
\end{build}

\begin{omsources}
\item A.~Moreira et T.~Muir, ``Volatility-managed portfolios'', \emph{Journal of Finance} 72(4), 2017.
\item S.~Cederburg, M.~S.~O'Doherty, F.~Wang et X.~S.~Yan, ``On the performance of volatility-managed portfolios'', \emph{Journal of Financial
Economics} 138(1), 2020.
\item C.~R.~Harvey, E.~Hoyle, R.~Korgaonkar, S.~Rattray, M.~Sargaison et O.~van Hemert, ``The impact of volatility targeting'', \emph{Journal of
Portfolio Management} 45(1), 2018.
\end{omsources}

\section{Exercices}

\begin{exercise}[$\star$]\label{exo:s1:volatility-targeting-and-risk-managed-portfolios:1}
Un fonds vise une volatilité de 10~\%. Quelle est son exposition lorsque la prévision est de 12~\%~? de 40~\%~?
\end{exercise}

\begin{exercise}[$\star$]\label{exo:s1:volatility-targeting-and-risk-managed-portfolios:2}
Des fonds détenant 0.2~\% de la valeur d'un marché réduisent leur exposition de 0.73 à 0.41 en une journée~; le marché échange 0.4~\% de sa valeur par jour. Quelle part du
volume de la journée vendent-ils~?
\end{exercise}

\begin{exercise}[$\star$]\label{exo:s1:volatility-targeting-and-risk-managed-portfolios:3}
Pourquoi la mise à l'échelle en variance inverse échange-t-elle davantage que la mise à l'échelle en volatilité inverse~?
\end{exercise}

\begin{exercise}[$\star\star$]\label{exo:s1:volatility-targeting-and-risk-managed-portfolios:4}
Le ratio de Sharpe vrai du marché synthétique est $0.06/0.16$. Quelle est sa valeur, quelle est l'erreur type d'une estimation sur dix ans, et celle du gain moyen sur
vingt échantillons lorsque les gains ont un écart-type de 0.22~?
\end{exercise}

\begin{exercise}[$\star\star$]\label{exo:s1:volatility-targeting-and-risk-managed-portfolios:5}
Pourquoi le ciblage améliore-t-il le ratio de Sharpe des actions et non celui des obligations dans les données de Harvey et de ses coauteurs~?
\end{exercise}

\begin{exercise}[$\star\star$]\label{exo:s1:volatility-targeting-and-risk-managed-portfolios:6}
Expliquer la boucle de rétroaction entre les ventes des fonds de ciblage et leurs propres prévisions de volatilité.
\end{exercise}

\begin{exercise}[$\star\star\star$]\label{exo:s1:volatility-targeting-and-risk-managed-portfolios:7}
\emph{Programmation.} Exécuter \texttt{spike\_flows} sur \texttt{base\_path()} avec une part de 0.5~\% et un impact de 0.3. Que se passe-t-il avant le pic, et pourquoi~?
\end{exercise}

\begin{exercise}[$\star\star\star$]\label{exo:s1:volatility-targeting-and-risk-managed-portfolios:8}
\emph{Trouver la faille.} «~Sur 1990--2020, le marché piloté par la volatilité a eu un alpha de 4~\% face au marché dans une régression d'écartement, donc nous allons le lancer.~»
\end{exercise}

\section{Problème~: vente mécanique}

\begin{problem}[{Problème du week-end --- ciblage et sa trace}]\label{pb:s1:volatility-targeting-and-risk-managed-portfolios:1}
Les marchés synthétiques du chapitre et les données publiques.

\medskip\noindent\textbf{Partie I --- Mécanique.}
\begin{enumerate}
\item Définir un \omterm{def:s1:volatility-targeting-and-risk-managed-portfolios:managed}{portefeuille piloté par la volatilité} et un \omterm{def:s1:volatility-targeting-and-risk-managed-portfolios:flow}{flux de rebalancement des produits à effet de levier}.
\item Énumérer les ingrédients d'une règle de \omterm{def:s1:trend-following:voltarget}{ciblage de volatilité}.
\item Quel est l'argument de Moreira et Muir~?
\item Que trouvent Cederburg et ses coauteurs~?
\end{enumerate}

\medskip\noindent\textbf{Partie II --- Ratios de Sharpe.}
\begin{enumerate}[resume]
\item Décrire l'expérience des vingt échantillons.
\item Donner les ratios de Sharpe moyens et le nombre d'échantillons améliorés.
\item Pourquoi un gain réel est-il invisible sur une seule décennie~?
\item Que trouvent Harvey et ses coauteurs, et que montre l'univers synthétique de contrats à terme~?
\end{enumerate}

\medskip\noindent\textbf{Partie III --- Flux.}
\begin{enumerate}[resume]
\item Décrire la simulation des flux et ses hypothèses.
\item Donner les ventes le lendemain de la chute et sur vingt jours.
\item Que fait la rétroaction à la chute~?
\item Pourquoi les ventes sont-elles prévisibles~?
\end{enumerate}

\medskip\noindent\textbf{Partie IV --- Le verdict.}
\begin{enumerate}[resume]
\item Énoncer le \emph{résultat nommé}~: l'amélioration du ratio de Sharpe grâce à la gestion de la volatilité et le flux agrégé, exprimé en fraction du volume après un pic.
\item Que fait le ciblage de façon fiable, sinon élever le ratio de Sharpe~?
\item Comment backtester honnêtement un \omterm{def:s1:volatility-targeting-and-risk-managed-portfolios:managed}{portefeuille piloté par la volatilité}~?
\item Quelle fiche de stratégie tire profit des autres~?
\item Quel lien ce chapitre entretient-il avec le ciblage de portefeuille du chapitre~19~?
\item De quoi auriez-vous besoin pour dimensionner le flux de désendettement sur un marché réel~?
\item Qu'est-ce qui rend le flux dangereux pour les fonds eux-mêmes~?
\item En une phrase~: que procure une cible de volatilité, et que coûte-t-elle au marché~?
\end{enumerate}
\end{problem}

\section{Questions d'entretien}

\begin{interviewq}[$\star$ \iqroles{chercheur}]\label{iq:s1:volatility-targeting-and-risk-managed-portfolios:1}
Pourquoi mettre à l'échelle l'exposition par l'inverse de la volatilité augmenterait-il un ratio de Sharpe~?
\end{interviewq}

\begin{interviewq}[$\star\star$ \iqroles{chercheur}]\label{iq:s1:volatility-targeting-and-risk-managed-portfolios:2}
Qu'y a-t-il de faux à juger un \omterm{def:s1:volatility-targeting-and-risk-managed-portfolios:managed}{portefeuille piloté par la volatilité} sur son alpha dans une régression d'écartement~?
\end{interviewq}

\begin{interviewq}[$\star\star$ \iqroles{trader}]\label{iq:s1:volatility-targeting-and-risk-managed-portfolios:3}
La volatilité a doublé aujourd'hui. Quels intervenants de marché doivent vendre demain, et comment estimeriez-vous le montant~?
\end{interviewq}

\begin{interviewq}[$\star\star$ \iqroles{risque}]\label{iq:s1:volatility-targeting-and-risk-managed-portfolios:4}
Quels risques une cible de volatilité réduit-elle, et lequel ajoute-t-elle~?
\end{interviewq}

\begin{interviewq}[$\star\star$ \iqroles{développeur}]\label{iq:s1:volatility-targeting-and-risk-managed-portfolios:5}
Concevez le processus quotidien qui fixe l'exposition d'un fonds à volatilité ciblée.
\end{interviewq}

\begin{interviewq}[$\star\star\star$ \iqroles{chercheur}]\label{iq:s1:volatility-targeting-and-risk-managed-portfolios:6}
Les rendements sont $r_t = \mu + \sigma_t\varepsilon_t$ avec un $\mu$ constant et $\sigma_t$ connu à l'avance. Montrez que la position $w_t = c/\sigma_t$ a un
ratio de Sharpe supérieur à celui d'une position constante, sauf si $\sigma_t$ est constant, et calculer le rapport des deux en fonction des moments de $1/\sigma_t$.
\end{interviewq}
